\documentclass[12pt,letterpaper]{article}



%%
%% Required packages:
%%
\usepackage{etex}
\usepackage{amsmath}
\usepackage{amssymb}
\usepackage{amstext}
\usepackage{amsthm}
\usepackage{fancyhdr,fancybox}
\usepackage{mathrsfs} % need for \mathscr command
\usepackage{multicol}
\usepackage[normalem]{ulem}
%\usepackage{pictex,pictexplus}
% \usepackage{pictexwd}
\usepackage{rotating}
\usepackage{url}
\usepackage{xfrac}
\usepackage{booktabs}
\usepackage{color,soul}
\usepackage{comment}

\usepackage{hyperref}
\hypersetup{
    colorlinks=true,     % false: boxed links; true: colored links
    linkcolor=blue,      % color of internal links
    citecolor=blue,      % color of links to bibliography
    filecolor=blue,      % color of file links
    urlcolor=blue        % color of external links
}


\usepackage{tikz}
\usetikzlibrary{backgrounds,calc}

%%
%% Common formatting commands
%%
\setlength{\parindent}{0in}
\setlength{\textwidth}{7in}
\setlength{\evensidemargin}{-0.25in}
\setlength{\oddsidemargin}{-0.25in}
\setlength{\parskip}{.5\baselineskip}
\setlength{\topmargin}{-0.5in}
\setlength{\textheight}{9in}

%               Problem and Part
\newcounter{problemnumber}
\newcounter{partnumber}
\newcounter{subpartnumber}
\newcommand{\Problem}{%
  \refstepcounter{problemnumber}%
  \setcounter{partnumber}{0}%
  \setcounter{subpartnumber}{0}%
  \item[\makebox{\hfil\textbf{\theproblemnumber.}\hfil}]%
  }
\newcommand{\InProblem}[2][1.6]{%
  \refstepcounter{problemnumber}%
  \setcounter{partnumber}{0}%
  \setcounter{subpartnumber}{0}%
  \textbf{\theproblemnumber.}\ \ \parbox[t]{#1in}{#2}%
}
\newcommand{\InSmallProblem}[1]{\InProblem[1.05]{#1}}
\newcommand{\NoProblem}[1]{%
  \mbox{\hfil\textbf{#1}\hfil}%
  }
\newcommand{\UnProblem}{%
  \refstepcounter{problemnumber}%
  \setcounter{partnumber}{0}%
  \setcounter{subpartnumber}{0}%
  \mbox{\hfil\textbf{\theproblemnumber.}\hfil}%
  }
\newcommand{\InFlexProblem}[1]{%
  \refstepcounter{problemnumber}%
  \setcounter{partnumber}{0}%
  \setcounter{subpartnumber}{0}%
  \textbf{\theproblemnumber.}\ \ #1%
}
%%  Part:
\newcommand{\Part}{%
  \renewcommand{\thepartnumber}{(\alph{partnumber})}%
  \refstepcounter{partnumber}
  \setcounter{subpartnumber}{0}%
  \item[(\alph{partnumber})]%
}
\newcommand{\InPart}[2][1.75]{%
  \renewcommand{\thepartnumber}{(\alph{partnumber})}%
  \refstepcounter{partnumber}%
  \setcounter{subpartnumber}{0}%
  (\alph{partnumber})\ \ \parbox[t]{#1in}{#2}%
}
\newcommand{\UnPart}{%
  \refstepcounter{partnumber}%
  \renewcommand{\thepartnumber}{(\alph{partnumber})}%
  \setcounter{subpartnumber}{0}%
  (\alph{partnumber})%
}
\newcommand{\InLargePart}[1]{\InPart[3.05]{#1}}
\newcommand{\InFullPart}[1]{\InPart[6.2]{#1}}
\newcommand{\InSmallPart}[1]{\InPart[1.05]{#1}}
%% SubPart:
\newcommand{\SubPart}{%
  \refstepcounter{subpartnumber}%
  \item[(\roman{subpartnumber})]%
  }
\newcommand{\InSubPart}[2][2.25]{%
  \stepcounter{subpartnumber}%
  (\roman{subpartnumber})\ \ \parbox[t]{#1in}{#2}%
}
\newcommand{\InSmallSubPart}[1]{\InSubPart[1.5]{#1}}
\newcommand{\InTinySubPart}[1]{%
  \stepcounter{subpartnumber}%
  (\roman{subpartnumber})\ \ \makebox{#1}%
}
\newcommand{\InMySubPart}[2][2.25]{%
  \stepcounter{subpartnumber}%
  \parbox[t]{#1in}{\makebox[0.3in][l]{(\roman{subpartnumber})}\ \ {#2}}%
}
\newcommand{\TrueFalse}{\textbf{T}\ \ \textbf{F}\ \ }
\newcommand{\TFPart}[1]{% 
  \stepcounter{partnumber}%
  \setcounter{subpartnumber}{0}%
  \item[(\alph{partnumber})] \TrueFalse\parbox[t]{5.5in}{#1}}

\newcommand{\Points}[1]{(#1 points) \ }
\newcommand{\Point}[1]{(#1 point) \ }


%% For Answers:
\newcounter{answernumber}
\newcommand{\Answer}{%
  \refstepcounter{answernumber}%
  \setcounter{partnumber}{0}%
  \setcounter{subpartnumber}{0}%
  \item[\makebox{\hfil\textbf{\theanswernumber.}\hfil}]%
  }


% Setting up the headers for the first page
%\chead{} % will default to what is typed in the file.
\fancyhead[L]{\textsc{Math 22a}}
\fancyhead[R]{\textsc{Fall 2024}}
\fancyfoot[R]{
	\vspace*{\fill}\footnotesize\textcopyright{} \emph{Harvard Math 22a}	
}
\renewcommand{\headrulewidth}{0pt}

%\fancyhead[C]{}
%	\chead{}	
%	\lhead{}
%	\rhead{}
%	\cfoot{}


% Putting copyright on the footer of each page
%\fancypagestyle{secondpage}{
%	\fancyhead[C]{TESTing again} % change this to be blank
%		%\chead{}	
%	\fancyhead[L]{bears} % change this to be blank
%	\fancyhead[R]{dogs} % change this to be blank
%	\fancyfoot[R]{ %keeps the footer the same
%		\vspace*{\fill}\footnotesize\textcopyright{} \emph{Harvard Math 22a}	
%	}
%}
%\renewcommand{\headrulewidth}{0pt}
%\pagestyle{fancy}
\pagestyle{empty}


%\lhead{}
%\rhead{}
%\rfoot{\vspace*{\fill}\footnotesize\textcopyright{} \emph{Harvard Math 22a}}
%\renewcommand{\headrulewidth}{0pt}
%\pagestyle{fancy}




%%
%% Common shortcut commands:
%%
\newcommand{\ds}{\displaystyle}
\newcommand{\sgn}{\operatorname{sgn}}
\renewcommand{\Pr}{\operatorname{P}}
\newcommand{\CPr}[3][{}]{\Pr_{\text{#1}}\left(#2\,|\,#3\right)}

\newcommand{\C}{\mathbb{C}}
\newcommand{\R}{\mathbb{R}}
\newcommand{\Z}{\mathbb{Z}}
\newcommand{\N}{\mathbb{N}}
\newcommand{\Q}{\mathbb{Q}}
\newcommand{\D}{\mathbb{D}}

\newcommand{\rref}{\operatorname{rref}}
\newcommand{\im}{\operatorname{im}}
\newcommand{\rank}{\operatorname{rank}}
\newcommand{\diag}{\operatorname{diag}}
\newcommand{\Span}{\operatorname{span}}
%\renewcommand{\vec}[1]{\mathbf{#1}}
\newcommand{\charpoly}{\mathscr{P}}
\newcommand{\nicefrac}[2]{\sfrac{#1}{#2}} % using xfrac package
\newcommand{\topology}{\mathcal{T}}
\newcommand{\Inn}{\operatorname{Inn}}

\newcommand{\blankbox}[2]{\fbox{\rule{#1}{0in}\rule{0in}{#2}}}

\newenvironment{myindentpar}[1]%
 {\begin{list}{}%
         {\setlength{\leftmargin}{#1}
          \setlength{\rightmargin}{\leftmargin}}%
         \item[]%
 }
 {\end{list}}
\newtheorem{theorem}{Theorem}
\newtheorem*{NStheorem}{Theorem}
\newtheorem{cor}[theorem]{Corollary}
\newtheorem*{claim}{Claim}
\newtheorem{defn}{Definition}

\newenvironment{amatrix}[1]{%
  \left(\begin{array}{@{}*{#1}{c}|c@{}}
}{%
  \end{array}\right)
}

%--------grstep
% For denoting a Gauss' reduction step.
% Use as: \grstep{\rho_1+\rho_3} or \grstep[2\rho_5 \\ 3\rho_6]{\rho_1+\rho_3}
\newcommand{\grstep}[2][\relax]{%
   \ensuremath{\mathrel{
       {\mathop{\longrightarrow}\limits^{#2\mathstrut}_{
                                     \begin{subarray}{l} #1 \end{subarray}}}}}}
\newcommand{\swap}{\leftrightarrow}




\makeatletter
\renewcommand*\env@matrix[1][*\c@MaxMatrixCols c]{%
	\hskip -\arraycolsep
	\let\@ifnextchar\new@ifnextchar
	\array{#1}}
\makeatother

\def\env@matrix{\hskip -\arraycolsep
	\let\@ifnextchar\new@ifnextchar
	\array{*\c@MaxMatrixCols c}}

\chead{\Large\textbf{Inverses}}% 

\begin{document}
\thispagestyle{fancy}

Recall from last class:

\begin{itemize}

\item Let $A=[\vec{a}_1 \; \dots \; \vec{a}_n]$ be an $m$ by $n$ matrix. We express $A$ by
$$A=\begin{bmatrix}
a_{11} & a_{12} & \hdots & a_{1n}\\ 
a_{21} & a_{22} & \hdots & a_{2n}\\ 
\vdots & & \ddots & \vdots\\
a_{n1} & a_{n2} & \hdots & a_{nn}\\ 
\end{bmatrix}$$
where $a_{ij}$ is the entry in the $i$-th row and $j$-th column.
\item {\bf Theorem:} Let $A$, $B$, and $C$ be matrices of the same size, and let $r$ and $s$ be real numbers. Then the following properties hold:
\begin{enumerate}
\item $A+B=B+A$.
\item $(A+B)+C=A+(B+C)$.
\item $A+0=A$.
\item $r(A+B)=rA+rB$.
\item $(r+s)A=rA+sA$.
\item $r(sA)=rsA$.
\end{enumerate}
\item If $A$ is a $m\times n$ matrix and $B$ is an $n\times p$ matrix with columns $\vec{b}_1, \dots, \vec{b}_p$, then $AB$ is the $m \times p$ matrix defined by $$AB=A[\vec{b}_1 \; \dots \; \vec{b}_p]=[A\vec{b}_1 \; \dots \; A\vec{b}_p].$$
\item Note $$(AB)_{ij}=a_{i1}b_{1j}+a_{i2}b_{2j}+\dots+a_{in}b_{nj}.$$
\item {\bf Theorem:}
\begin{enumerate}
\item $A(BC)=(AB)C$.
\item $A(B+C)=AB+AC$.
\item $(B+C)A=BA+CA$.
\item $r(AB)=(rA)B=A(rB)$.
\item $I_m A =A = AI_n$.
\end{enumerate}
\item Given an $m\times n$ matrix $A$, then the {\bf transpose} of $A$, denoted $A^{T}$ is the $n \times m$ matrix given by taking the columns of $A$ as rows.
\item {\bf Theorem:} Let $A$ be an $m$ by $n$ matrix, $B$ be an $n$ by $p$ matrix, and $r\in \R$.
\begin{enumerate}
\item $(A^{T})^{T}=A.$
\item $(A+B)^{T}=A^{T}+B^{T}.$
\item $(rA)^{T}=r A^{T}$.
\item $(AB)^{T}=B^{T} A^{T}$.
\end{enumerate}
\end{itemize}

\newpage

\begin{defn}
An $n$ by $n$ matrix $A$ is {\bf invertible} if there is an $n$ by $n$ matrix $C$ so that $AC=I_n$ and $CA=I_n$. We denote the inverse of $A$ by $A^{-1}$.
\end{defn}

{\bf Natural Questions:}
\begin{enumerate}
\item Do inverses always exist?
\item If not, when do they exist?
\item If they exist, are they unique?
\item How can we find them?
\end{enumerate}

\begin{enumerate}

\Problem Let $A=\begin{bmatrix} a & b\\ c & d\\ \end{bmatrix}$. Determine when $A$ is invertible and find $A^{-1}$. 

\newpage

{\bf Theorem:} Let $A=\begin{bmatrix} a & b \\ c & d\\ \end{bmatrix}$. If $ad-bc\neq 0$, then $A$ is invertible  and $$A^{-1}= \frac{1}{ad-bc} \begin{bmatrix} d & -b\\ -c &a \\ \end{bmatrix}.$$

\begin{defn}
We define the {\bf determinant} of $A=\begin{bmatrix} a & b \\ c & d\\ \end{bmatrix}$, denoted by $\text{det}(A)$, by $$\text{det}(A)=ad-bc.$$
\end{defn}

{\bf Theorem:} If $A$ is an $n$ by $n$ invertible matrix, then for each $\vec{b}\in \R^n$, $A\vec{x}=\vec{b}$ has a unique solution; namely $\vec{x}=A^{-1}\vec{b}$.

\Problem Prove the Theorem.

\vfill

\newpage

{\bf Theorem:} Suppose $A$ and $B$ are $n$ by $n$ invertible matrices, then the following hold. 
\begin{enumerate}
\item $A^{-1}$ is invertible and $(A^{-1})^{-1}=A$.
\item $(AB)^{-1}=B^{-1} A^{-1}$.
\item $(A^{T})^{-1}=(A^{-1})^{T}$.
\end{enumerate}

\Problem Prove the Theorem.

\vfill

\newpage

\begin{defn}
An {\bf elementary matrix} $E$ is a matrix obtained from the identity matrix by a single elementary row operation.
\end{defn}

\Problem Give some examples of elementary matrices.

\vfill

{\bf Theorem:} $A$ is invertible if and only if $A$ is row equivalent to the identity matrix. Furthermore, any sequence of elementary row operations that reduces $A$ to $I_n$ also transforms $I_n$ into $A^{-1}$.

\Problem Prove the theorem.

\vfill
\vfill
\vfill

\newpage

\Problem 
\Part Let $A=\begin{bmatrix}
1 & -2 & -1\\
-1 &5 & 6\\
5 &-4 & 5\\
\end{bmatrix}$. Find $A^{-1}$ if possible.

\vfill

\Part Let $A=\begin{bmatrix}
1 & 0 & -2\\
-3 &1 & 4\\
2 &-3 & 4\\
\end{bmatrix}$. Find $A^{-1}$ if possible.

\vfill

\newpage

{\bf Invertible Matrix Theorem} Let $A$ be an $n \times n$ matrix, then the following are all equivalent.
\begin{enumerate}
\item $A$ is invertible.
\item $A$ is row-equivalent to $I_n$.
\item $A$ has $n$ pivot positions.
\item $A\vec{x}=\vec{0}$ has only the trivial solution.
\item The columns of $A$ are linearly independent.
\item $x\mapsto A\vec{x}$ is one-to-one (injective).
\item For all $\vec{b}\in \R^n$, $A\vec{x}=\vec{b}$ is consistent.
\item The columns of $A$ span $\R^n$.
\item $\vec{x} \mapsto A\vec{x}$ is onto (surjective).
\item There exists a matrix $C$ such that $CA=I_n$.
\item There exists a matrix $D$ such that $AD=I_n$.
\item $A^{T}$ is invertible.
\end{enumerate}

\Problem Prove the theorem.

\newpage

\Problem Determine if the following matrices are invertible.
\Part Let $A=\begin{bmatrix}
3 & 0 & 0\\
-3 &-4 & 0\\
8 & 5 & -3\\
\end{bmatrix}$.

\vfill

\Part Let $A=\begin{bmatrix}
1 &2 & 3\\
4 &5 & 6\\
7 & 8 & 9\\
\end{bmatrix}$.

\vfill

\newpage

{\bf Theorem:} Let $T: \R^n \to \R^n$ be linear and let $A$ be the associated matrix. Then $T$ is invertible if and only if $A$ is invertible.

\Problem Prove the theorem.
 

\end{enumerate}

\newpage
\chead{\Large\textbf{Inverses -- Answers / Solutions}}%
\lhead{}
\rhead{}
\thispagestyle{fancy}

\begin{enumerate}
\Answer We want to find $e$, $f$, $g$ and $h$ such that 
$$\begin{bmatrix}
a & b\\
c & d\\
\end{bmatrix}\begin{bmatrix}
e & f\\
g &h\\
\end{bmatrix}=\begin{bmatrix}
1 & 0\\
0 & 1\\
\end{bmatrix}.$$
Thus we get the following system of equations:
\begin{align*}
	ae+gb&=1\\
	af+bh&=0\\
	ce+dg&=0\\
	cf+dh&=1\\
\end{align*}

Multplying the first equation by $c$ and the third equation by $-a$ and adding them together yields: $bcg-adg=c$, and assuming $ad-bc\neq0$, we get $g=\frac{-c}{ad-bc}$.

Continuing in a similar manner we get $e=\frac{d}{ad-bc}$, $h=\frac{a}{ad-bc}$, and $f=\frac{-b}{ad-bc}$.

Thus the matrix is invertible if and only if $ad-bc\neq0$, and we have an explicit formula for the inverse matrix.

\Answer Since $A$ is invertible, we know $A^{-1}$ exists. Let $\vec{x}=A^{-1}\vec{b}$, then we check $\vec{x}$ is a solution. Here $A\vec{x}=AA^{-1}\vec{b}=\vec{b}$, and hence $\vec{x}$ is a solution.

Now we check it is the unique solution. Suppose $\vec{u}$ is another solution, then $A\vec{u}=\vec{b}$. Since $A$ is invertible, we can multiply both sides by $A^{-1}$ to get $A^{-1}A\vec{u}=A^{-1}\vec{b}$, and hence $\vec{u}=A^{-1}\vec{b}=\vec{x}$. Thus $\vec{x}$ is the only solution.

\Answer 

\begin{enumerate}
	
	\Part Suppose $A$ is invertible. We check that $A$ is the inverse of $A^{-1}$; namely, $AA^{-1}=I_n$ and $A^{-1}A=I_n$. Thus the inverse of $A^{-1}$ is $A$; in symbols $(A^{-1})^{-1}=A$.
	
	\Part Again suppose $A$ and $B$ are both invertible matrices. We check that $B^{-1}A^{-1}$ is the inverse of $AB$: $ABB^{-1}A^{-1}=I_n$ and $B^{-1}A^{-1}AB=I_n$. Thus by the definition of inverse matrix, $B^{-1}A^{-1}$ is the inverse of $AB$; i.e. $(AB)^{-1}=B^{-1}A^{-1}$.
	
	\Part Suppose $A$ is an invertible matrix. We want to prove that the inverse of $A^{T}$ is $(A^{-1})^{T}$. Again we check the conditions in the definition of matrix inverse: $A^{T}(A^{-1})^{T}=(A^{-1}A)^{T}=(I_n)^{T}=I_n$ and $(A^{-1})^{T}A^{T}=(AA^{-1})^{T}=(I_n)^{T}=I_n$ where we used $(AB)^{T}=B^{T}A^{T}$. Thus $(A^{T})^{-1}=(A^{-1})^{T}$.
	
\end{enumerate}
  
  \Answer Let $$E_1=\begin{bmatrix} 2 & 0 \\ 0 & 1\\ \end{bmatrix}, E_2=\begin{bmatrix} 0 & 1\\ 1 &0\\ \end{bmatrix}, E_3=\begin{bmatrix}
  1 &0\\ 2 & 1\\ \end{bmatrix}.$$ Now $E_1$ corresponds to scaling the first row by 2, $E_2$ corresponds to interchanging the first two rows, and $E_3$ corresponds to adding two times the first row to the second row. Notice each elementary matrix is invertible by undoing the corresponding elementary row operation. Thus $$E_1^{-1}=\begin{bmatrix} \frac{1}{2} & 0 \\ 0 & 1\\ \end{bmatrix}, E_2^{-1}=\begin{bmatrix} 0 & 1\\ 1 &0\\ \end{bmatrix}, E_3^{-1}=\begin{bmatrix}
  1 &0\\ -2 & 1\\ \end{bmatrix}.$$
  
  \newpage
  
  {\bf Note:}
  \begin{itemize}
  	\item Elementary matrices are invertible.
  	\item If $E$ is an elementary matrix and $A$ is a matrix, then $EA$ corresponds to performing $E$'s row operation on $A$.
  \end{itemize}
  
  \Answer 
  
  \begin{proof} We prove both implications directly. Let $A$ be an $n$ by $n$ matrix.
  	
  $\implies$ Suppose $A$ is invertible. We previously proved that for all $\vec{b}\in \mathbb{R}^n$, the matrix equation $A\vec{x}=\vec{b}$ is consistent. We proved in chapter one that the matrix equation $A\vec{x}=\vec{b}$ is consistent if and only if $A$ has a pivot position in each row. Since $A$ is an $n$ by $n$ matrix, we know there is also a pivot position in each column. Thus the reduced row echelon form of $A$ is $I_n$, and hence $A$ is row equivalent $I_n$.
  
  $\impliedby$ Suppose $A$ is row equivalent to $I_n$. Thus there exists elementary matrices $E_1, \dots, E_p$ such that $$A\sim E_1 A \sim E_2 E_1A \sim \dots \sim E_p E_{p-1}\dots E_2 E_1 A=I_n.$$ 
  
  Since elementary matrices are inertible and the product of invertible matrices is invertible, we get $A$ is invertible and $A^{-1}=E_p\dots E_2 E_1$. Thus, if we apply the same sequence of row operations that transforms $A$ to $I_n$ to the identity matrix we get $E_p\dots E_2 E_1 I_n=A^{-1}$.
  \end{proof}
  
  \Answer
  
 \begin{enumerate} 
  \Part Let $A=\begin{bmatrix}
  1 & -2 & -1\\
  -1 &5 & 6\\
  5 &-4 & 5\\
  \end{bmatrix}$. We try augment $A$ by $I_n$ and try to use the result that if $A$ is invertible, then $[A|I_n]\sim [I_n | A^{-1}]$. Thus
  
  \begin{align*}
  \begin{pmatrix}[ccc|ccc]
  1 & -2 & -1 & 1 & 0 & 0\\
  -1 &5 & 6 & 0 & 1 & 0\\
  5 &-4 & 5 & 0 & 0 &1\\
  \end{pmatrix}& \sim  \begin{pmatrix}[ccc|ccc]
  1 & -2 & -1 & 1 & 0 & 0\\
  0 &3 & 5 & 1 & 1 & 0\\
  0 & 6 & 10 & -5 & 0 &1\\
\end{pmatrix} \sim \begin{pmatrix}[ccc|ccc]
1 & -2 & -1 & 1 & 0 & 0\\
0 &3 & 5 & 1 & 1 & 0\\
0 &0 & 0 & -7 & -2 &1\\
\end{pmatrix},
  \end{align*}
and now we're stuck. Since $A$ is not row-equivalent to $I_n$ (which we're seeing on left), the algorithm fails to find an inverse for $A$. Thus we know $A$ is not invertible.

\Part Let  $A=\begin{bmatrix}
1 & 0 & -2\\
-3 &1 & 4\\
2 &-3 & 4\\
\end{bmatrix}$. We again augment $A$ by $I_n$ and try to use the result that if $A$ is invertible, then $[A|I_n]\sim [I_n | A^{-1}]$. Thus

\begin{align*}
	\begin{pmatrix}[ccc|ccc]
		1 & 0 & -2 & 1 & 0 & 0\\
		-3 &1 & 4 & 0 & 1 & 0\\
		2 &-3 & 4 & 0 & 0 &1\\
	\end{pmatrix}& \sim  \begin{pmatrix}[ccc|ccc]
		1 & 0 & -2 & 1 & 0 & 0\\
		0 &1 & -2 & 3 & 1 & 0\\
		0 & -3 & 8 & -2 & 0 &1\\
	\end{pmatrix}\\ &\sim \begin{pmatrix}[ccc|ccc]
		1 & 0 & -2 & 1 & 0 & 0\\
		0 &1 & -2 & 3 & 1 & 0\\
		0 &0 & 2 & 7 & 3 &1\\
	\end{pmatrix}\\ &\sim
\begin{pmatrix}[ccc|ccc]
	1 & 0 & 0 & 8 & 3 &1 \\
	0 &1 & 0 & 10 & 4 & 1\\
	0 &0 & 2 & \frac{7}{2} & \frac{3}{2} &\frac{1}{2}\\
\end{pmatrix},
\end{align*}
and hence $A^{-1}=\begin{bmatrix}
8 & 3 & 1\\
10 & 4 & 1\\
\frac{7}{2} & \frac{3}{2} & \frac{1}{2}\\
\end{bmatrix}.$

\end{enumerate}
  
  \Answer 
  
 \begin{proof} There are many ways to organize all the implications. Here I try to use previous results from class as much as possible.
 
 $(a)\implies (j)$ by definition.
 
 $(j) \implies (d)$ by Theorem 2.5 (or above on this handout).
 
 $(d) \implies (c)$ by Theorem 1.4.
 
 $(c) \implies (b)$ since $A$ is square.
 
 $(b) \implies (a)$ by Theorem 2.7 (or above on this handout).
 
 Thus we have established $(a) \iff (b) \iff (c) \iff (d) \iff (j)$.
 
Note $(a)\implies (k)$ by definition of invertibility. We prove $(k) \implies (g)$. Suppose $D$ is an $n\times n$ matrix such that $AD=I_n$. Let $\vec{b} \in \mathbb{R}^n$, and define $\vec{y} =D\vec{b}$. Now $A\vec{y}=AD\vec{b} =I_n\vec{b} =\vec{b}$, and hence $\vec{y}$ is a solution to the matrix equation $A\vec{x}=\vec{b}$. Thus the matrix equation is always consistent, and hence $(g)$ follows. Now we prove $(g)\implies (a)$. By Theorem 1.4, $A\vec{x}=\vec{b}$ being consistent for any $\vec{b}$ is equivalent to $A$ having a pivot position in each row. Since $A$ is square, we know $A$ must be row equivalent to the identity matrix. By the theorem above, $A$ is invertible, and hence $(a)$ holds. 

Finally note $(d) \iff (e) \iff (f)$ and $(g) \iff (h) \iff (i)$ follows from Theorem 1.12. Finally we prove $(a)\iff (l)$ above on this handout.
 \end{proof}
  
  \Answer 
  
  \begin{enumerate}
  \Part We could row-reduce $A$ to determine if $A$ is row equivalent to the identity matrix, but in this case we can see the columns of $A$ are linearly independent, and hence we know $A$ must be invertible from the invertible matrix theorem.
  
  \Part We have previously seen this matrix does not span all of $\mathbb{R}^3$, and hence we can use the invertible matrix theorem to conclude $A$ is not invertible.
  
  \end{enumerate}
  
  \Answer
  
  \begin{proof} We prove both implications directly from the definitions.
  
  $\implies$ Suppose $T$ is invertible, then there exists a function $T^{-1} : \mathbb{R}^n \to \mathbb{R}^n$ such that $T(T^{-1}(\vec{x}))=\vec{x}$ and $T^{-1}(T(\vec{x}))=\vec{x}$ for all $\vec{x} \in \mathbb{R}^n$. If we knew that $T^{-1}$ was linear, then it would have an associated matrix we could use.
  
  {\bf Claim:} $T^{-1}$ is linear.
  
  {\bf Proof of Claim:} Let $\vec{x}, \vec{y} \in \mathbb{R}^n$ and $c \in \mathbb{R}$. Since $T$ is invertible, it is bijective. Thus there exists unique $\vec{a}, \vec{b} \in \mathbb{R}^n$ such that $T(\vec{a})=\vec{x}$ and $T(\vec{b})=\vec{y}$. Now we check the conditions for linearity: $$T^{-1}(\vec{x}+\vec{y})=T^{-1}(T(\vec{a})+T(\vec{b}))=T^{-1}(T(\vec{a}+\vec{b})=\vec{a}+\vec{b}=T^{-1}(\vec{x})+T^{-1}(\vec{y}).$$ and $$T^{-1}(c\vec{x})=T^{-1}(c T(\vec{a}))=T^{-1}(T(c\vec{a}))=c \vec{a} = c T^{-1}(\vec{x}).$$ Thus $T^{-1}$ is linear.
  
  Since $T^{-1}$ is linear, it has an associated matrix $D$. Now $$AD\vec{x} =T(T^{-1}(\vec{x}))=\vec{x}$$ and $$DA\vec{x} =T^{-1}(T(\vec{x})) =\vec{x}.$$ Thus $D$ is the inverse matrix of $A$, and hence $A$ is invertible. 
  
  $\impliedby$ Suppose $A$ is invertible. Define $S: \mathbb{R}^n \to \mathbb{R}^n$ by $S(\vec{x})=A^{-1}\vec{x}$. Then $T(S(\vec{x}))=AA^{-1}\vec{x} = \vec{x}$ and $S(T(\vec{x})) =A^{-1}A\vec{x}=\vec{x}$. Thus $S$ is the inverse function of $T$, and hence $T$ is invertible. 
  \end{proof}
\end{enumerate}

\end{document}


%%% Local ables: 
%%% mode: latex
%%% TeX-master: t
%%% End: 
